% --- INSTITUTIONAL TEMPLATE FED v7.4.9 ---
% Estándar de Reporteo APA 7 (Babel-Free)
\documentclass[12pt,spanish]{article}

\usepackage{graphicx}
\usepackage{geometry}
\usepackage{indentfirst} 
\usepackage{caption}
\usepackage{floatrow} 

% Forzar estilos de caption antes de cualquier carga
\floatsetup[figure]{capposition=top}
\floatsetup[table]{capposition=top}

\usepackage{booktabs}
\usepackage{float}
\usepackage{longtable}
\usepackage{threeparttable}
\usepackage{array}
\usepackage{calc}
\usepackage{amsmath, amssymb}
\usepackage{fancyhdr}
\usepackage{mathptmx} 
\usepackage[T1]{fontenc}
\usepackage{setspace} 
\usepackage{ragged2e} 
\usepackage{titlesec}
\usepackage{url}
\def\UrlBreaks{\do\/\do-\do\.\do\_\do\?\do\=\do\&\do\%}
\usepackage[hidelinks,breaklinks]{hyperref}

% --- Pandoc Compatibility Block (Iron Clad) ---
\providecommand{\tightlist}{%
  \setlength{\itemsep}{0pt}\setlength{\parskip}{0pt}}

\usepackage{xcolor}
\usepackage{fancyvrb}
\usepackage{framed}

% --- Code Highlighting (Pandoc Standard) ---
\AtBeginDocument{
  \definecolor{shadecolor}{RGB}{248,248,248}
}

% Compatibilidad Pandoc 3.x
\ifdefined\pandocbounded\else
  \newcommand{\pandocbounded}[1]{#1}
\fi

% --- Pandoc Citeproc (CSL) compatibility ---
\newlength{\cslhangindent}
\setlength{\cslhangindent}{1.27cm} 
\newlength{\csllabelwidth}
\setlength{\csllabelwidth}{3em}
\newcommand{\citeproctext}{} 
\newcommand{\citeprocitem}[2]{\item} 
\newenvironment{CSLReferences}[2] 
 {\begin{list}{}{
  \setlength{\leftmargin}{\cslhangindent}
  \setlength{\parsep}{0pt}
  \setlength{\itemsep}{6pt}
  \ifodd #1
   \setlength{\leftmargin}{\leftmargin + \csllabelwidth}
   \setlength{\itemindent}{-\csllabelwidth}
  \fi
  }
  \renewcommand{\bibitem}[2][]{\item} % Elimina los corchetes [ ]
 }
 {\end{list}}
\newcommand{\CSLBlock}[1]{#1\hfill\break}
\newcommand{\CSLLeftMargin}[1]{\parbox[t]{\csllabelwidth}{#1}}
\newcommand{\CSLRightInline}[1]{\parbox[t]{\linewidth - \csllabelwidth}{#1}\break}
\newcommand{\CSLIndent}[1]{\hspace{\cslhangindent}#1}

\makeatletter
\@ifundefined{paragraph}{\newcounter{paragraph}}{}
\@ifundefined{subparagraph}{\newcounter{subparagraph}}{}
\makeatother

% Solución al error 'No counter none defined' en longtable
\makeatletter
\newcounter{none}
\@ifpackageloaded{caption}{
  \DeclareCaptionType{none}
  \captionsetup[none]{name=Tabla}
}{}
\makeatother

\hypersetup{
  colorlinks=true,
  linkcolor=blue,
  filecolor=blue,
  citecolor=blue,
  urlcolor=blue
}

% --- Macros Institucionales ---
\newcommand{\fuente}[1]{
  \vspace{4pt}
  \begin{flushleft}
    \small
    \textit{Nota.} #1
  \end{flushleft}
}

% --- Localización Manual Determinista (Sin Babel) ---
\AtBeginDocument{
  \renewcommand{\tablename}{Tabla}
  \renewcommand{\figurename}{Figura}
  \renewcommand{\contentsname}{Índice}
  \renewcommand{\listtablename}{Índice de Tablas}
  \renewcommand{\listfigurename}{Índice de Figuras}
  \renewcommand{\abstractname}{Resumen}
  \renewcommand{\refname}{Referencias}
  
  \RaggedRight
  \setlength{\parindent}{1.27cm}
  \setlength{\parskip}{0pt}
}

\DeclareCaptionLabelSeparator{newline}{\\ \vspace{2pt}}
\captionsetup{
  position=top,
  justification=RaggedRight,
  singlelinecheck=false,
  labelfont=bf,
  textfont=it,
  labelsep=newline,
  font=small,
  skip=10pt
}
\captionsetup[figure]{name=Figura}
\captionsetup[table]{name=Tabla}

% Márgenes APA 7
\geometry{
  a4paper,
  margin=25.4mm,
  headheight=15mm
}

% --- Encabezados ---
\pagestyle{fancy}
\fancyhf{} 
\fancyhead[L]{\includegraphics[height=1.2cm]{\logoUNL}}
\fancyhead[R]{\includegraphics[height=1.2cm]{\logoECONOMIA}}
\renewcommand{\headrulewidth}{0pt}
\fancyfoot[R]{\thepage}

\fancypagestyle{plain}{
  \fancyhf{}
  \fancyhead[L]{\includegraphics[height=1.2cm]{\logoUNL}}
  \fancyhead[R]{\includegraphics[height=1.2cm]{\logoECONOMIA}}
  \fancyfoot[R]{\thepage}
  \renewcommand{\headrulewidth}{0pt}
}

% --- Títulos ---
\titleformat{\section}{\normalfont\normalsize\bfseries\centering}{\thesection}{1em}{}
\titleformat{\subsection}{\normalfont\normalsize\bfseries}{\thesubsection}{1em}{}
\titleformat{\subsubsection}{\normalfont\normalsize\bfseries\itshape}{\thesubsubsection}{1em}{}
\titlespacing*{\section}{0pt}{12pt}{6pt}
\titlespacing*{\subsection}{0pt}{12pt}{6pt}
\titlespacing*{\subsubsection}{0pt}{12pt}{6pt}

\newenvironment{referenciasAPA}{
  \setlength{\parindent}{-1.27cm}
  \setlength{\leftskip}{1.27cm}
  \setlength{\parskip}{6pt}
}{\par}

\usepackage{fontspec}
\newcommand{\logoUNL}{/home/erick-fcs/Capital_Workstation/capital-workstation-libs/src/ecs_quantitative/reporting/assets/logos/logo_unl.png}
\newcommand{\logoECONOMIA}{/home/erick-fcs/Capital_Workstation/capital-workstation-libs/src/ecs_quantitative/reporting/assets/logos/logo_economia.jpg}
\renewcommand{\RaggedRight}{\justifying}\makeatletter\@ifpackageloaded{caption}{}{\usepackage{caption}}\AtBeginDocument{%
\ifdefined\contentsname
  \renewcommand*\contentsname{Table of contents}
\else
  \newcommand\contentsname{Table of contents}
\fi
\ifdefined\listfigurename
  \renewcommand*\listfigurename{List of Figures}
\else
  \newcommand\listfigurename{List of Figures}
\fi
\ifdefined\listtablename
  \renewcommand*\listtablename{List of Tables}
\else
  \newcommand\listtablename{List of Tables}
\fi
\ifdefined\figurename
  \renewcommand*\figurename{Figure}
\else
  \newcommand\figurename{Figure}
\fi
\ifdefined\tablename
  \renewcommand*\tablename{Table}
\else
  \newcommand\tablename{Table}
\fi
}\@ifpackageloaded{float}{}{\usepackage{float}}
\floatstyle{ruled}
\@ifundefined{c@chapter}{\newfloat{codelisting}{h}{lop}}{\newfloat{codelisting}{h}{lop}[chapter]}
\floatname{codelisting}{Listing}\newcommand*\listoflistings{\listof{codelisting}{List of Listings}}\makeatother\makeatletter\makeatother\makeatletter\@ifpackageloaded{caption}{}{\usepackage{caption}}
\@ifpackageloaded{subcaption}{}{\usepackage{subcaption}}\makeatother

\begin{document}

\begin{center}
    {\bfseries \Large UNIVERSIDAD NACIONAL DE LOJA}\\[0.3cm]
    {\bfseries \large FACULTAD JURÍDICA, SOCIAL Y
ADMINISTRATIVA}\\[0.2cm]
    {\bfseries CARRERA DE ECONOMÍA}\\[1.5cm]
    
    {\bfseries \large POLÍTICA ECONÓMICA}\\[1cm]
    
    {\bfseries \Large Tarea de Aprendizaje Colaborativo: Análisis
Cuantitativo de una Política Pública en Ecuador}\\[1.5cm]
\end{center}

\begin{center}
    \setstretch{1.1}
                        Luis Romero \\
                    Luis Abad \\
                    Dayana Francelina Bustamante Benitez \\
                    Erick Fabricio Condoy Seraquive \\
                \vspace{0.5cm}
    Econ. Maria Gabriela Moreno Hurtado \\ 
    22 de junio de 2026
\end{center}

\vspace{1cm}

\doublespacing
\setlength{\parindent}{1.27cm}

\newpage

\section*{Resumen Ejecutivo}\label{resumen-ejecutivo}
\addcontentsline{toc}{section}{Resumen Ejecutivo}

La presente evaluación analiza la efectividad y sostenibilidad de los
incentivos a la generación de empleo contemplados en la Ley Orgánica de
Eficiencia Económica y Generación de Empleo (Registro Oficial Suplemento
N.º 461, 2023), política impulsada por el gobierno del presidente Daniel
Noboa con el objetivo explícito de fomentar el empleo, incrementar la
recaudación tributaria y estimular la inversión. La normativa introduce
deducciones adicionales del 50\% y 75\% sobre sueldos de nuevas plazas
laborales, exoneraciones fiscales en Zonas Francas y un régimen de
remisión tributaria. Los datos utilizados provienen de la Encuesta
Nacional de Empleo, Desempleo y Subempleo (ENEMDU) del Instituto
Nacional de Estadística y Censos (INEC), el Banco Central del Ecuador
(BCE), el Servicio de Rentas Internas (SRI) y la legislación analizada.
La metodología empleada es de carácter descriptivo-comparativo, centrada
en el análisis de indicadores laborales y fiscales entre los periodos
2023 y 2024. Los resultados evidencian un incremento notable en la
recaudación tributaria, la cual alcanzó los USD 20.131 millones en 2024,
lo que representa un aumento del 15,6\%. No obstante, se observa un
deterioro simultáneo en el mercado laboral: la tasa de empleo adecuado
disminuyó del 36,3\% en 2023 al 35,9\% en 2024, mientras que la
informalidad ascendió al 58\%, cifra que representa el nivel más alto
registrado desde 2007. Asimismo, la inversión extranjera directa se
contrajo a USD 232 millones, alcanzando su mínimo en catorce años. Se
concluye que la política fue parcialmente efectiva en su componente
fiscal, pero ineficaz en su dimensión laboral, mientras que su
sostenibilidad resulta cuestionable debido a que los ingresos
extraordinarios derivados de la remisión poseen un carácter no
recurrente.

\section{Introducción}\label{introducciuxf3n}

Ecuador enfrenta un mercado laboral estructuralmente precario. Al cierre
de 2023, más de la mitad de la población ocupada se desempeñaba en
condiciones de informalidad, con empleos sin contrato, sin acceso a
seguridad social y con ingresos inferiores al salario básico unificado
de USD 450 mensuales. En ese contexto, el gobierno del presidente Daniel
Noboa promulgó la Ley Orgánica de Eficiencia Económica y Generación de
Empleo (en adelante, LEEGE) el 20 de diciembre de 2023, una norma
calificada como de urgencia económica, que introduce reformas
tributarias estructurales, crea el régimen de Zonas Francas, y establece
incentivos fiscales directos para la generación de nuevas plazas de
trabajo (Asamblea Nacional del Ecuador, 2023).

El problema que esta política busca resolver es la insuficiencia de
empleo adecuado y la alta tasa de informalidad laboral, fenómenos que
limitan la productividad, la recaudación del IESS y la calidad de vida
de la población ecuatoriana. El instrumento principal es un incentivo
tributario: deducciones adicionales al Impuesto a la Renta para
empleadores que generen nuevas plazas de trabajo, especialmente para
jóvenes entre 18 y 29 años, personas con obligaciones alimenticias y ex
privados de libertad (Art. 6, LEEGE, 2023).

La pregunta de evaluación central es: ¿Los incentivos a la generación de
empleo establecidos en la LEEGE fueron efectivos para reducir la
informalidad y ampliar el empleo adecuado en Ecuador durante 2024, y qué
tan sostenible es esta política en términos fiscales, económicos e
institucionales? El informe se organiza en once secciones siguiendo la
estructura metodológica requerida: descripción de la política, lógica de
intervención, marco conceptual, datos e indicadores, estadística
descriptiva, estrategia metodológica, resultados, sostenibilidad y
conclusiones.

\section{Descripción de la Política Evaluada y Lógica de
Intervención}\label{descripciuxf3n-de-la-poluxedtica-evaluada-y-luxf3gica-de-intervenciuxf3n}

\subsection{Caracterización de la
política}\label{caracterizaciuxf3n-de-la-poluxedtica}

La LEEGE fue aprobada por la Asamblea Nacional el 19 de diciembre de
2023, sancionada por el Ejecutivo ese mismo día y publicada en el
Registro Oficial Suplemento N.° 461 el 20 de diciembre de 2023. Es una
Ley Orgánica de urgencia económica ---modalidad constitucional que
impone plazos acelerados de tratamiento legislativo (Art. 137 de la
Constitución)--- y se encuentra en vigencia con última modificación al
11 de enero de 2024 (Asamblea Nacional del Ecuador, 2023).

La institución responsable de su implementación es el Servicio de Rentas
Internas (SRI) en el componente tributario, el Comité Estratégico de
Promoción y Atracción de Inversiones (CEPAI) en Zonas Francas, y el
Ministerio del Trabajo en el registro de nuevos contratos. El objetivo
explícito declarado en el Art. 1 de la ley es: ``el impulso inmediato
del empleo, el incremento de la recaudación tributaria y el incentivo de
la inversión'' (LEEGE, Art. 1, 2023). La población objetivo son los
empleadores (personas naturales y sociedades) del sector formal que
generen nuevas plazas de trabajo con aporte al IESS.

Los instrumentos centrales para la generación de empleo son: (a)
deducción adicional del 50\% sobre sueldos y salarios de nuevas plazas
para jóvenes de 18-29 años y personas con obligaciones alimenticias; (b)
deducción del 75\% si los nuevos empleados son jóvenes graduados de
universidades o institutos públicos; (c) deducción del 75\% para nuevas
plazas en los sectores de construcción y agricultura; y (d) deducción
del 75\% para empleadores que contraten a ex privados de libertad con
sentencia superior a un año (LEEGE, Art. 6, numerales 9.2 y 9.3, 2023).
Adicionalmente, la ley crea el régimen de Zonas Francas con 0\% de IR
por cinco años, exoneración de IVA e ISD, y un periodo de remisión
tributaria del 100\% de intereses, multas y recargos para obligaciones
generadas hasta el 31 de diciembre de 2023 (LEEGE, Disposición
Transitoria Primera, 2023).

\subsection{Lógica de intervención de la
política}\label{luxf3gica-de-intervenciuxf3n-de-la-poluxedtica}

Para evaluar la cadena de impacto de los incentivos, se sistematiza la
lógica de intervención de la LEEGE en la Tabla 1:

\textbf{Tabla 1}\\
\emph{Matriz de Lógica de Intervención de la LEEGE}

{\def\LTcaptype{none} % do not increment counter
\begin{longtable}[]{@{}
  >{\raggedright\arraybackslash}p{(\linewidth - 2\tabcolsep) * \real{0.1186}}
  >{\raggedright\arraybackslash}p{(\linewidth - 2\tabcolsep) * \real{0.8814}}@{}}
\toprule\noalign{}
\begin{minipage}[b]{\linewidth}\raggedright
Elemento
\end{minipage} & \begin{minipage}[b]{\linewidth}\raggedright
Contenido
\end{minipage} \\
\midrule\noalign{}
\endhead
\bottomrule\noalign{}
\endlastfoot
\textbf{Problema inicial} & Alta informalidad laboral
(\textasciitilde55,7\% al cierre de 2023) y baja tasa de empleo adecuado
(36,3\%). \\
\textbf{Instrumento de política} & Deducciones adicionales del IR (50\%
y 75\%) por nuevas plazas de trabajo registradas en el Ministerio del
Trabajo con aporte al IESS. \\
\textbf{Canal de transmisión} & La reducción del costo tributario neto
del trabajo formal incentiva a los empleadores a formalizar relaciones
laborales o crear nuevas plazas en el sector formal. \\
\textbf{Resultado esperado} & Aumento de la tasa de empleo adecuado,
reducción de la informalidad y ampliación de la base de afiliados al
IESS. \\
\textbf{Indicador de evaluación} & Tasa de empleo adecuado
(ENEMDU-INEC), tasa de informalidad y número de nuevos afiliados al
IESS. \\
\textbf{Supuesto clave} & Los empleadores del sector formal son
sensibles al costo tributario del trabajo y ampliarán la nómina formal
ante reducciones de ese costo. \\
\textbf{Riesgo o efecto no deseado} & La deducción solo aplica a
contribuyentes del régimen general de IR, excluyendo a las microempresas
que concentran la mayor parte del empleo informal; el incentivo puede
financiar rotación de personal en lugar de empleo neto nuevo. \\
\end{longtable}
}

\emph{Nota.} Elaboración propia con base en LEEGE (2023) y datos
ENEMDU-INEC (2023).

\section{Marco Conceptual de
Evaluación}\label{marco-conceptual-de-evaluaciuxf3n}

La política analizada se inscribe dentro del objetivo de pleno empleo
definido por Musgrave (1959) como uno de los tres objetivos primordiales
de la política fiscal, junto con la estabilidad de precios y el
crecimiento económico. En el marco de los instrumentos de política
económica, los incentivos tributarios al empleo corresponden a lo que
Tinbergen (1952) clasifica como instrumentos indirectos: actúan sobre
las decisiones privadas sin imponer mandatos directos, modificando las
señales de precio relativo del trabajo formal frente al informal.

La relación entre fiscalidad y mercado laboral ha sido extensamente
analizada. Saez et al.~(2019) demuestran, en el contexto de los
subsidios salariales estadounidenses (EITC), que la elasticidad del
empleo formal frente a créditos tributarios es significativa pero
concentrada en los márgenes extensivos (decisión de participar) más que
intensivos (horas trabajadas). En economías con alta informalidad, Levy
(2008) advierte que los incentivos tributarios al empleo formal pueden
ser capturados por el sector ya formalizado, generando ineficiencia
distributiva sin reducción neta de la informalidad.

En Ecuador, la dualidad formal-informal responde a factores
estructurales ---baja productividad de microempresas, costo del IESS,
rigideces del salario mínimo--- que los incentivos tributarios
difícilmente atacan de raíz (Banco Mundial, 2022). El conflicto entre
objetivos es claro: la deducción adicional reduce la recaudación de
corto plazo por IR (efecto negativo sobre el objetivo fiscal) mientras
que pretende ampliar la base imponible del IESS (efecto positivo de
largo plazo). La sostenibilidad fiscal depende entonces de que el
aumento de afiliados al IESS y de nuevos contribuyentes formales
compense la renuncia fiscal inducida por la deducción adicional, lo que
requiere evaluación cuantitativa.

\section{Datos e Indicadores}\label{datos-e-indicadores}

Los datos utilizados en esta evaluación provienen exclusivamente de
fuentes oficiales. La principal fuente de información laboral es la
Encuesta Nacional de Empleo, Desempleo y Subempleo (ENEMDU), producida
mensual y anualmente por el Instituto Nacional de Estadística y Censos
(INEC). Para los indicadores fiscales y macroeconómicos se utilizan los
boletines del Banco Central del Ecuador (BCE) y los reportes de
recaudación del Servicio de Rentas Internas (SRI). El texto legal de
referencia es la LEEGE publicada en el Registro Oficial Suplemento N.°
461.

\textbf{Tabla 2}\\
\emph{Indicadores de evaluación de la política}

{\def\LTcaptype{none} % do not increment counter
\begin{longtable}[]{@{}
  >{\raggedright\arraybackslash}p{(\linewidth - 8\tabcolsep) * \real{0.1333}}
  >{\raggedright\arraybackslash}p{(\linewidth - 8\tabcolsep) * \real{0.3167}}
  >{\raggedright\arraybackslash}p{(\linewidth - 8\tabcolsep) * \real{0.1333}}
  >{\raggedright\arraybackslash}p{(\linewidth - 8\tabcolsep) * \real{0.1000}}
  >{\raggedright\arraybackslash}p{(\linewidth - 8\tabcolsep) * \real{0.3167}}@{}}
\toprule\noalign{}
\begin{minipage}[b]{\linewidth}\raggedright
Indicador
\end{minipage} & \begin{minipage}[b]{\linewidth}\raggedright
Definición
\end{minipage} & \begin{minipage}[b]{\linewidth}\raggedright
Fuente
\end{minipage} & \begin{minipage}[b]{\linewidth}\raggedright
Periodo
\end{minipage} & \begin{minipage}[b]{\linewidth}\raggedright
Relación con la política
\end{minipage} \\
\midrule\noalign{}
\endhead
\bottomrule\noalign{}
\endlastfoot
Tasa de empleo adecuado & \% de la PEA con trabajo de al menos
40h/semana e ingresos \(\ge\) SBU & ENEMDU-INEC (anual) & 2021-2024 &
Resultado directo esperado de la política \\
Tasa de informalidad & \% de ocupados sin RUC en establecimientos
\textless100 trabajadores & ENEMDU-INEC & 2023-2024 & Resultado directo;
la ley busca reducirla \\
Tasa de desempleo & \% de la PEA sin trabajo y en búsqueda activa &
ENEMDU-INEC (anual) & 2021-2024 & Indicador de contexto laboral \\
Recaudación tributaria total & Total recaudado por el SRI (USD millones)
& SRI Boletín 2024 & 2023-2024 & Efectividad del componente fiscal de la
ley \\
Recaudación IR & Impuesto a la Renta recaudado (USD millones) & SRI &
2023-2024 & Mide si la deducción redujo el IR neto \\
Remisión tributaria & Monto recaudado por remisión de intereses y multas
& SRI, 2024 & 2024 & Ingreso extraordinario no recurrente \\
IED (USD millones) & Inversión Extranjera Directa neta & BCE & 2023-2024
& Sostenibilidad: ¿la ley atrajo inversión? \\
PIB real (variación \%) & Crecimiento anual del PIB en volumen & BCE,
Cuentas Nacionales & 2023-2024 & Contexto macroeconómico general \\
\end{longtable}
}

\emph{Nota.} SBU = Salario Básico Unificado. PEA = Población
Económicamente Activa. Elaboración propia.

\section{Estadística Descriptiva}\label{estaduxedstica-descriptiva}

\subsection{Mercado laboral: tendencia
2021-2024}\label{mercado-laboral-tendencia-2021-2024}

La Tabla 3 presenta la evolución de los principales indicadores
laborales en Ecuador, según la ENEMDU anual del INEC. La tasa de empleo
adecuado mostró una recuperación sostenida desde 2021 (32,5\%) hasta
2023 (36,3\%), impulsada por la normalización post-pandémica. Sin
embargo, en 2024 ---primer año completo de vigencia de la LEEGE--- la
tasa cayó ligeramente a 35,9\%, revirtiéndose la tendencia de mejora
(INEC, 2024).

\textbf{Tabla 3}\\
\emph{Evolución de indicadores laborales en Ecuador (2021-2024)}

{\def\LTcaptype{none} % do not increment counter
\begin{longtable}[]{@{}lccccc@{}}
\toprule\noalign{}
Indicador (\% PEA) & 2021 & 2022 & 2023 & 2024 & Variación 2023-2024 \\
\midrule\noalign{}
\endhead
\bottomrule\noalign{}
\endlastfoot
Tasa de empleo adecuado & 32,5 & 34,4 & 36,3 & 35,9 & -0,4 pp \\
Tasa de subempleo & 23,2 & 22,2 & 19,6 & 21,0 & +1,4 pp \\
Tasa de desempleo & 5,2 & 4,4 & 3,8 & 3,7 & -0,1 pp \\
Tasa de informalidad (dic.) & n.d. & n.d. & 55,7 & 58,0 & +2,3 pp \\
\end{longtable}
}

\emph{Nota.} pp = puntos porcentuales. n.d. = no disponible con
desagregación comparable. Fuente: INEC, ENEMDU anual (2024). Elaboración
propia.

La reducción del empleo adecuado en 2024 se produjo a pesar de la
entrada en vigor de la ley en enero de dicho año. Resulta
particularmente crítico que la informalidad alcanzara el 58\% en
diciembre de 2024, lo cual representa el nivel más alto registrado desde
2007 (INEC, 2024; La Hora, 2025). Esta tendencia sugiere que los
incentivos tributarios de la LEEGE resultaron insuficientes para
contrarrestar factores estructurales y coyunturales adversos, tales como
la recesión económica, la crisis energética, la inseguridad y el
incremento del Impuesto al Valor Agregado (IVA) al 15\% dispuesto en la
misma normativa.

\subsection{Indicadores fiscales: recaudación tributaria
2023-2024}\label{indicadores-fiscales-recaudaciuxf3n-tributaria-2023-2024}

En contraste con los resultados del componente laboral, el desempeño
fiscal derivado de la aplicación de la ley fue significativamente
positivo. La recaudación total del SRI ascendió a USD 20.131 millones en
2024, lo que supone un incremento del 15,6\% respecto a los USD 17.420
millones recaudados en 2023 (SRI, 2025). Este crecimiento se explica
fundamentalmente por tres fuentes extraordinarias: en primer lugar, la
remisión tributaria, que generó un aporte de USD 562 millones
provenientes de contribuyentes en mora; en segundo lugar, las
autorretenciones de los 520 grandes contribuyentes, mecanismo
establecido por la LEEGE; y, finalmente, las contribuciones temporales
aplicadas al sistema financiero (SRI, 2025; Primicias, 2024).

\textbf{Tabla 4}\\
\emph{Comparativa de indicadores fiscales y de inversión (2023-2024)}

{\def\LTcaptype{none} % do not increment counter
\begin{longtable}[]{@{}lccc@{}}
\toprule\noalign{}
Concepto & 2023 (USD mill.) & 2024 (USD mill.) & Variación (\%) \\
\midrule\noalign{}
\endhead
\bottomrule\noalign{}
\endlastfoot
Recaudación tributaria total & 17.420 & 20.131 & +15,6\% \\
Impuesto a la Renta (IR) & n.d. & 6.639 & n.d. \\
IVA & n.d. & 9.767 & n.d. \\
Remisión tributaria (extraordinario) & --- & 562 & --- \\
IED (flujo anual) & ≈ 530* & 232 & -56,2\% \\
\end{longtable}
}

\emph{Nota.} *Estimación BCE. IED = Inversión Extranjera Directa.
Fuente: SRI (2025); BCE (2024); Primicias (2025). Elaboración propia.

La Inversión Extranjera Directa (IED), indicador fundamental para medir
la capacidad de la normativa en la atracción de capitales productivos,
sufrió una contracción severa al situarse en USD 232 millones en 2024,
cifra que representa el nivel más bajo en catorce años (Primicias,
2025). Este hecho evidencia que el marco normativo favorable instaurado
por la LEEGE no fue suficiente para mitigar el impacto de un entorno
adverso caracterizado por la inseguridad, crisis eléctricas prolongadas
y una contracción del Producto Interno Bruto (PIB) estimada en el 0,9\%
para el ejercicio fiscal 2024 (BCE, 2024b).

\section{Estrategia Metodológica de
Evaluación}\label{estrategia-metodoluxf3gica-de-evaluaciuxf3n}

Se adopta una estrategia descriptiva-comparativa de tipo antes/después
(pre-post), basada en la comparación de indicadores laborales y fiscales
entre el año previo a la ley (2023) y el primer año completo de vigencia
(2024). Esta elección metodológica es la más adecuada dado que:

\begin{enumerate}
\def\labelenumi{\alph{enumi})}
\tightlist
\item
  La política se aplicó de manera universal a todo el territorio
  nacional, lo que impide construir un grupo de control geográfico o
  sectorial.
\item
  Los datos disponibles no permiten identificar el número exacto de
  plazas creadas mediante las deducciones adicionales (no existe
  publicación administrativa del SRI con ese desagregado).
\item
  El periodo de vigencia es de un año, insuficiente para aplicar métodos
  de series de tiempo con adecuada potencia estadística.
\end{enumerate}

El análisis descriptivo examina tendencias, tasas de variación y
comparaciones antes/después de los indicadores seleccionados. Se
reconoce explícitamente que este diseño no permite inferir causalidad:
los cambios observados en 2024 son el resultado conjunto de múltiples
factores (ley, recesión, crisis energética, inseguridad), sin que sea
posible aislar el efecto atribuible exclusivamente a los incentivos
laborales de la LEEGE. La evaluación se limita a establecer asociaciones
y contrastar la hipótesis de mejora laboral con la evidencia disponible.

\section{Estrategia Empírica y Justificación
Metodológica}\label{estrategia-empuxedrica-y-justificaciuxf3n-metodoluxf3gica}

No se aplica un modelo econométrico formal. La justificación es la
siguiente: dado que la LEEGE fue implementada simultáneamente en todo el
territorio nacional y no existe un grupo no tratado, no es posible
construir un estimador de diferencias en diferencias (DiD) válido. La
ausencia de datos administrativos publicados sobre el número de
contribuyentes que efectivamente accedieron a la deducción adicional por
nuevas plazas (numerales 9.2 y 9.3, Art. 6 de la LEEGE) impide
cuantificar el tratamiento. Un modelo de Mínimos Cuadrados Ordinarios
(MCO) de corte transversal o de series de tiempo con un año de
observación post-intervención carecería de potencia estadística y
produciría estimaciones sesgadas e ininterpretables.

En este contexto, la evaluación descriptiva comparativa, bien
justificada y con indicadores pertinentes, es más sólida
epistemológicamente que un modelo econométrico mal identificado. Como
señala Angrist y Pischke (2009), la credibilidad de una evaluación
depende de la calidad del diseño de identificación, no de la
sofisticación del modelo. El análisis se concentra en tres preguntas
evaluativas operacionales:

\begin{enumerate}
\def\labelenumi{\arabic{enumi})}
\tightlist
\item
  ¿Cambió la tasa de empleo adecuado en la dirección esperada?
\item
  ¿Se redujo la informalidad?
\item
  ¿Fueron los resultados fiscales atribuibles a los mecanismos de la
  ley?
\end{enumerate}

\section{Resultados y Discusión de
Efectividad}\label{resultados-y-discusiuxf3n-de-efectividad}

\subsection{Efectividad laboral}\label{efectividad-laboral}

Los hallazgos no corroboran la hipótesis de efectividad en el componente
laboral. La tasa de empleo adecuado disminuyó del 36,3\% en 2023 al
35,9\% en 2024 (INEC, 2024), lo cual representa una pérdida de 0,4
puntos porcentuales. Al analizar los datos mensuales, el deterioro se
torna más evidente: en diciembre de 2024 la tasa descendió al 33\%, en
comparación con el 35,9\% registrado en diciembre de 2023, lo que
implica una reducción de 2,9 puntos porcentuales (Bloomberg Línea,
2025). Simultáneamente, la informalidad se incrementó del 55,7\% al
58,0\%, lo cual significa que el indicador contrario al objetivo de la
ley se degradó de manera considerable y alcanzó su nivel más
desfavorable desde 2007.

En noviembre de 2024, el número de personas con empleo adecuado fue de
2.905.637, frente a 3.037.803 en noviembre de 2023, una pérdida de más
de 100.000 empleos plenos en un año (El Comercio, 2024). El efecto fue
heterogéneo: el empleo adecuado urbano fue del 44,5\% mientras el rural
apenas alcanzó el 19,4\% (INEC, 2024). La brecha de género persiste:
41,4\% para hombres versus 28,4\% para mujeres. Estos resultados
sugieren que el impacto de la ley fue nulo o negativo en términos de
empleo formal durante 2024.

\subsection{Efectividad fiscal}\label{efectividad-fiscal}

El componente fiscal presenta resultados positivos, aunque con
importantes matices sobre su atribución y sostenibilidad. La recaudación
del SRI creció 15,6\% en 2024, superando la meta inicial. Sin embargo,
el SRI reconoce que ``buena parte de este aumento se sustentó en
ingresos extraordinarios como la contribución temporal del sistema
financiero y la remisión tributaria aprobada en la Ley de Eficiencia
Económica, que por sí sola aportó más de USD 561 millones'' (La Hora,
2025). Esto implica que el crecimiento recurrente de la recaudación es
significativamente menor al reportado en el agregado.

El mecanismo de autorretención de los 520 Grandes Contribuyentes, creado
por la LEEGE, fue el principal impulsor del Impuesto a la Renta, que
alcanzó USD 6.639 millones en 2024 (SRI, 2025). Este resultado es
positivo y recurrente. Sin embargo, la deducción adicional para nuevas
plazas de trabajo implica menor recaudación neta de IR para los
empleadores que accedan a ella, generando un trade-off fiscal que no ha
sido cuantificado públicamente por el SRI.

\textbf{Tabla 5}\\
\emph{Resumen de resultados de efectividad de la política}

{\def\LTcaptype{none} % do not increment counter
\begin{longtable}[]{@{}
  >{\raggedright\arraybackslash}p{(\linewidth - 6\tabcolsep) * \real{0.3750}}
  >{\raggedright\arraybackslash}p{(\linewidth - 6\tabcolsep) * \real{0.3750}}
  >{\centering\arraybackslash}p{(\linewidth - 6\tabcolsep) * \real{0.1250}}
  >{\centering\arraybackslash}p{(\linewidth - 6\tabcolsep) * \real{0.1250}}@{}}
\toprule\noalign{}
\begin{minipage}[b]{\linewidth}\raggedright
Variable
\end{minipage} & \begin{minipage}[b]{\linewidth}\raggedright
Resultado observado
\end{minipage} & \begin{minipage}[b]{\linewidth}\centering
Dirección esperada
\end{minipage} & \begin{minipage}[b]{\linewidth}\centering
¿Efectivo?
\end{minipage} \\
\midrule\noalign{}
\endhead
\bottomrule\noalign{}
\endlastfoot
Empleo adecuado (2023 \(\rightarrow\) 2024) & 36,3\% \(\rightarrow\)
35,9\% (-0,4 pp) & Aumento & No \\
Informalidad (dic. 2023 \(\rightarrow\) dic. 2024) & 55,7\%
\(\rightarrow\) 58,0\% (+2,3 pp) & Reducción & No \\
Recaudación tributaria & USD 17.420 M \(\rightarrow\) 20.131 M (+15,6\%)
& Aumento & Sí (parcial) \\
IED (2023 \(\rightarrow\) 2024) & \(\approx\) USD 530 M \(\rightarrow\)
232 M (-56,2\%) & Aumento & No \\
Remisión tributaria & USD 562 millones (no recurrente) & Ingreso
adicional & Transitorio \\
\end{longtable}
}

\emph{Nota.} pp = puntos porcentuales. Fuente: INEC (2024); SRI (2025);
BCE (2024); Primicias (2025). Elaboración propia.

\section{Evaluación de
Sostenibilidad}\label{evaluaciuxf3n-de-sostenibilidad}

\textbf{Tabla 6}\\
\emph{Evaluación de sostenibilidad de la política económica}

{\def\LTcaptype{none} % do not increment counter
\begin{longtable}[]{@{}
  >{\raggedright\arraybackslash}p{(\linewidth - 4\tabcolsep) * \real{0.0205}}
  >{\raggedright\arraybackslash}p{(\linewidth - 4\tabcolsep) * \real{0.9538}}
  >{\centering\arraybackslash}p{(\linewidth - 4\tabcolsep) * \real{0.0256}}@{}}
\toprule\noalign{}
\begin{minipage}[b]{\linewidth}\raggedright
Dimensión
\end{minipage} & \begin{minipage}[b]{\linewidth}\raggedright
Evaluación
\end{minipage} & \begin{minipage}[b]{\linewidth}\centering
Clasificación
\end{minipage} \\
\midrule\noalign{}
\endhead
\bottomrule\noalign{}
\endlastfoot
Económica & El deterioro del mercado laboral en 2024 (empleo adecuado
-0,4 pp, informalidad +2,3 pp) sugiere que los beneficios esperados no
se materializaron. La caída del PIB al 0,9\% en 2024 contextualizó
adversamente la política. & Baja \\
Fiscal & La remisión tributaria (USD 562 M) es no recurrente. El
crecimiento de la recaudación depende de medidas extraordinarias que se
agotan en 2024. La deducción adicional por plazas nuevas reduce el IR
recurrente. & Baja-media \\
Social & La informalidad alcanzó el 58\% ---máximo histórico reciente---
y la pobreza subió del 26\% al 28\%. Los beneficios laborales no
llegaron a los grupos más vulnerables (microempresas, trabajadores
informales). & Baja \\
Institucional & El CEPAI, el SRI y el Ministerio del Trabajo cuentan con
capacidades para administrar la norma. El régimen de Zonas Francas
requiere infraestructura institucional aún en construcción. & Media \\
Ambiental & La ley excluye actividades extractivas y contaminantes de
las Zonas Francas (Art. 50, LEEGE, 2023). No hay efectos ambientales
directos negativos identificados en el primer año. & Media-alta \\
Política & La ley fue sancionada bajo la figura de urgencia económica,
limitando el debate legislativo. Esto genera riesgo de cuestionamiento
político-judicial posterior. & Media-baja \\
\end{longtable}
}

\emph{Nota.} Elaboración propia con base en LEEGE (2023); INEC (2024);
SRI (2025); BCE (2024).

En conclusión, la política se clasifica como efectiva en su componente
fiscal de corto plazo, aunque esta efectividad es parcial y se sustenta
en ingresos de carácter no recurrente. No obstante, la política resulta
ineficaz en su objetivo laboral primordial. La sostenibilidad general se
considera baja en las dimensiones económica, fiscal y social, lo que
sitúa a esta normativa en la categoría de política efectiva únicamente
en el corto plazo y carente de sostenibilidad en los términos evaluados.

\section{Conclusiones y
Recomendaciones}\label{conclusiones-y-recomendaciones}

\subsection{Conclusiones}\label{conclusiones}

La evaluación de la LEEGE revela una dicotomía crítica entre los
objetivos fiscales y laborales del Estado. Mientras que la normativa
logró una consolidación fiscal inmediata mediante mecanismos agresivos
de recaudación como la autorretención y la remisión, este éxito
financiero no se tradujo en una dinamización del mercado de trabajo. La
brecha observada sugiere que los incentivos tributarios, por sí solos,
resultan inoperantes cuando el entorno macroeconómico está marcado por
la recesión y la inseguridad. La política pública, por tanto, ha
priorizado la liquidez fiscal sobre la estabilidad social del empleo
pleno, exacerbando la precariedad laboral e informalidad.

En términos de implicaciones, la dependencia de ingresos no recurrentes
para sostener el presupuesto estatal compromete la sostenibilidad fiscal
a largo plazo. La erosión del empleo adecuado indica que la reducción
del costo impositivo del trabajo es insuficiente para compensar los
altos costos operativos y el riesgo país que enfrentan las empresas.
Para rectificar este rumbo, se requiere una transición desde una
política recaudatoria reactiva hacia una estrategia integral de fomento
productivo que aborde las barreras estructurales de la inversión y la
productividad.

\subsection{Recomendaciones}\label{recomendaciones}

Se presentan las siguientes recomendaciones estratégicas:

\begin{enumerate}
\def\labelenumi{\arabic{enumi})}
\tightlist
\item
  \textbf{Reforma de incentivos:} Reorientar las deducciones adicionales
  exclusivamente hacia las microempresas y PYMES, asegurando que el
  beneficio esté condicionado a la formalización verificable de
  trabajadores previamente informales.
\item
  \textbf{Transparencia y monitoreo:} Implementar un sistema de reporte
  trimestral por parte del SRI que detalle el uso de beneficios
  tributarios y su correlación directa con la creación neta de puestos
  de trabajo por sector.
\item
  \textbf{Protección social proactiva:} Evaluar la sustitución de
  deducciones por subsidios directos a la seguridad social en sectores
  de baja productividad, mitigando el riesgo de informalidad
  estructural.
\end{enumerate}

\section*{Referencias}\label{referencias}
\addcontentsline{toc}{section}{Referencias}

\begingroup
\setlength{\parindent}{-1.27cm}
\setlength{\leftskip}{1.27cm}
\setlength{\parskip}{6pt}

Angrist, J. D., \& Pischke, J. S. (2009). \emph{Mostly harmless
econometrics: An empiricist's companion}. Princeton University Press.

Asamblea Nacional del Ecuador. (2023). Ley Orgánica de Eficiencia
Económica y Generación de Empleo. \emph{Registro Oficial Suplemento N.°
461}. https://www.lexis.com.ec

Banco Central del Ecuador. (2024a). \emph{Informe de resultados cuentas
nacionales trimestrales: Primer trimestre 2024}.
https://contenido.bce.fin.ec/documentos/informacioneconomica/cuentasnacionales/trimestrales/Informe\_CNTITrim2024.pdf

Banco Central del Ecuador. (2024b). \emph{Estadísticas del sector real:
Crecimiento anual del PIB}.
https://contenido.bce.fin.ec/documentos/informacioneconomica/SectorReal/ix\_SectorRealPrin.html

Banco Central del Ecuador. (2025). \emph{La economía ecuatoriana creció
3,4\% en el primer trimestre de 2025}.
https://www.bce.fin.ec/la-economia-ecuatoriana-crecio-34-en-el-primer-trimestre-de-2025/

Banco Mundial. (2022). \emph{Ecuador: Diagnóstico sistemático de país}.
Grupo del Banco Mundial.

Bloomberg Línea. (2025, 26 de enero). El empleo adecuado en Ecuador cae
y la informalidad alcanza récord.
https://www.bloomberglinea.com/latinoamerica/ecuador/el-empleo-adecuado-en-ecuador-cae-y-la-informalidad-alcanza-record/

CISS Bienestar. (2024, 13 de noviembre). Más del 50\% de trabajadores
tiene un empleo informal en Ecuador.
https://ciss-bienestar.org/2024/11/14/mas-del-50-de-trabajadores-tiene-un-informal-en-ecuador/

El Comercio. (2024, 23 de diciembre). Empleo adecuado se reduce en
Ecuador: Mire las últimas cifras del 2024.
https://www.elcomercio.com/actualidad/negocios/empleo-adecuado-se-reduce-en-ecuador-mire-las-ultimas-cifras-del-2024/

Expreso. (2025, 27 de enero). La recesión de 2024 se tradujo en más
pobreza e informalidad.
https://www.expreso.ec/actualidad/economia/recesion-2024-tradujo-pobreza-e-informalidad-229431.html

Instituto Nacional de Estadística y Censos. (2024). \emph{Encuesta
Nacional de Empleo, Desempleo y Subempleo: ENEMDU anual 2024}.
https://www.ecuadorencifras.gob.ec/documentos/web-inec/EMPLEO/2024/anual/Boletin\_tecnico\_anual\_enero-diciembre\_2024.pdf

La Hora. (2025, 21 de abril). La mitad de los ecuatorianos con empleo
cerraron 2024 con ingresos de USD 354,6 al mes o menos.
https://www.lahora.com.ec/archivo/La-mitad-de-los-ecuatorianos-con-empleo-cerraron-2024-con-ingresos-de-3546-al-mes-o-menos-2025

Levy, S. (2008). \emph{Good intentions, bad outcomes: Social policy,
informality, and economic growth in Mexico}. Brookings Institution
Press.

Lexis Ecuador. (2025, 6 de enero). Inversión extranjera directa en
Ecuador cae más del 50\% en el tercer trimestre de 2024.
https://www.lexis.com.ec/noticias/inversion-extranjera-directa-en-ecuador-cae-mas-del-50-en-el-tercer-trimestre-de-2024

Musgrave, R. A. (1959). \emph{The theory of public finance: A study in
public economy}. McGraw-Hill.

Primicias. (2024a, 25 de agosto). Remisión tributaria dejó USD 562
millones en recaudación de impuestos.
https://www.primicias.ec/economia/remision-tributaria-dejo-usd-562-millones-recaudacion-impuestos-77245/

Primicias. (2024b, 21 de abril). Menos trabajo pleno y más desempleo:
Las cifras laborales de marzo de 2024.
https://www.primicias.ec/noticias/economia/empleo-desempleo-ecuador-trabajo-inec-referendum-informalidad/

Primicias. (2025, 6 de abril). La inversión extranjera de 2024 en
Ecuador fue la más baja en 14 años.
https://www.primicias.ec/economia/inversion-extranjera-ecuador-china-estados-unidos-93425/

Primicias. (2024c, 26 de diciembre). Estas son las cifras de empleo con
las que Ecuador cierra el 2024.
https://www.primicias.ec/economia/cifras-empleo-ecuador-desempleo-trabajo-2024-inec-86360/

Saez, E., Schoefer, B., \& Seim, D. (2019). Payroll taxes, firm
behavior, and rent sharing: Evidence from a young workers' tax cut in
Sweden. \emph{American Economic Review}, \emph{109}(5), 1717--1763.
https://doi.org/10.1257/aer.20171937

Servicio de Rentas Internas. (2025). \emph{Boletín N.° NAC-COM-25-003:
Las ventas crecieron 2,2\% e impulsaron la recaudación tributaria del
año 2024}. SRI. https://www.sri.gob.ec

Tinbergen, J. (1952). \emph{On the theory of economic policy}.
North-Holland Publishing.

\endgroup

\section*{Anexos}\label{anexos}
\addcontentsline{toc}{section}{Anexos}

\subsection{Anexo A. Diccionario de
variables}\label{anexo-a.-diccionario-de-variables}

\textbf{Tabla 7}\\
\emph{Diccionario de variables analizadas}

{\def\LTcaptype{none} % do not increment counter
\begin{longtable}[]{@{}
  >{\raggedright\arraybackslash}p{(\linewidth - 6\tabcolsep) * \real{0.0833}}
  >{\raggedright\arraybackslash}p{(\linewidth - 6\tabcolsep) * \real{0.7500}}
  >{\raggedright\arraybackslash}p{(\linewidth - 6\tabcolsep) * \real{0.0833}}
  >{\raggedright\arraybackslash}p{(\linewidth - 6\tabcolsep) * \real{0.0833}}@{}}
\toprule\noalign{}
\begin{minipage}[b]{\linewidth}\raggedright
Variable
\end{minipage} & \begin{minipage}[b]{\linewidth}\raggedright
Definición operacional
\end{minipage} & \begin{minipage}[b]{\linewidth}\raggedright
Fuente
\end{minipage} & \begin{minipage}[b]{\linewidth}\raggedright
Frecuencia
\end{minipage} \\
\midrule\noalign{}
\endhead
\bottomrule\noalign{}
\endlastfoot
Tasa de empleo adecuado & \% de la PEA que trabaja \(\ge\) 40 h/semana
con ingresos \(\ge\) SBU & ENEMDU-INEC & Mensual / Anual \\
Tasa de informalidad & \% de ocupados en establecimientos \textless{}
100 trab. sin RUC & ENEMDU-INEC & Mensual \\
Tasa de desempleo & \% de la PEA sin empleo y en búsqueda activa &
ENEMDU-INEC & Mensual / Anual \\
Recaudación SRI (USD M) & Total de tributos administrados por el SRI &
SRI & Mensual / Anual \\
IED (USD M) & Inversión Extranjera Directa neta registrada por el BCE &
BCE & Trimestral / Anual \\
PIB real (\%) & Variación porcentual del PIB en volumen (base 2018=100)
& BCE - CNT & Trimestral / Anual \\
SBU & Salario Básico Unificado: USD 450 (2023) / USD 460 (2024) &
Ministerio del Trabajo & Anual \\
PEA & Población Económicamente Activa (\(\ge\) 15 años) & ENEMDU-INEC &
Continua \\
\end{longtable}
}

\emph{Nota.} Elaboración propia.

\subsection{Anexo B. Artículos clave de la LEEGE referenciados en el
informe}\label{anexo-b.-artuxedculos-clave-de-la-leege-referenciados-en-el-informe}

\textbf{Tabla 8}\\
\emph{Artículos clave de la LEEGE referenciados}

{\def\LTcaptype{none} % do not increment counter
\begin{longtable}[]{@{}
  >{\raggedright\arraybackslash}p{(\linewidth - 2\tabcolsep) * \real{0.1667}}
  >{\raggedright\arraybackslash}p{(\linewidth - 2\tabcolsep) * \real{0.8333}}@{}}
\toprule\noalign{}
\begin{minipage}[b]{\linewidth}\raggedright
Artículo
\end{minipage} & \begin{minipage}[b]{\linewidth}\raggedright
Contenido relevante
\end{minipage} \\
\midrule\noalign{}
\endhead
\bottomrule\noalign{}
\endlastfoot
Art. 1 (Objeto) & Impulso inmediato del empleo, incremento de
recaudación tributaria e incentivo a la inversión. \\
Art. 6, num. 9.2 & Deducción adicional del 50\% sobre sueldos de nuevas
plazas para jóvenes 18-29 años; 75\% si son egresados de instituciones
públicas. \\
Art. 6, num. 9.3 & Deducción adicional del 75\% por nuevas plazas para
ex privados de libertad con sentencia \textgreater{} 1 año. \\
Art. 34 (Zonas Francas) & Definición y régimen especial aduanero,
tributario y financiero de las Zonas Francas. \\
Art. 35 (Objetivos ZF) & Las Zonas Francas promueven generación de
empleo, atracción de inversión y competitividad nacional. \\
Art. 50.3 (Régimen tributario ZF) & 0\% de IR por 5 años; luego 15\%
fijo. Exoneración de IVA, ISD y tributos al comercio exterior. \\
Disp. Transitoria Primera & Remisión del 100\% de intereses, multas y
recargos por obligaciones hasta dic. 2023, con pago hasta jul. 2024. \\
\end{longtable}
}

\emph{Nota.} Elaboración propia.

\end{document}
